% Options for packages loaded elsewhere
\PassOptionsToPackage{unicode}{hyperref}
\PassOptionsToPackage{hyphens}{url}
%
\documentclass[
  12pt,
  a4paper,
  DIV=11,
  numbers=noendperiod]{scrreprt}

\usepackage{amsmath,amssymb}
\usepackage{setspace}
\usepackage{iftex}
\ifPDFTeX
  \usepackage[T1]{fontenc}
  \usepackage[utf8]{inputenc}
  \usepackage{textcomp} % provide euro and other symbols
\else % if luatex or xetex
  \usepackage{unicode-math}
  \defaultfontfeatures{Scale=MatchLowercase}
  \defaultfontfeatures[\rmfamily]{Ligatures=TeX,Scale=1}
\fi
\usepackage{lmodern}
\ifPDFTeX\else  
    % xetex/luatex font selection
\fi
% Use upquote if available, for straight quotes in verbatim environments
\IfFileExists{upquote.sty}{\usepackage{upquote}}{}
\IfFileExists{microtype.sty}{% use microtype if available
  \usepackage[]{microtype}
  \UseMicrotypeSet[protrusion]{basicmath} % disable protrusion for tt fonts
}{}
\usepackage{xcolor}
\setlength{\emergencystretch}{3em} % prevent overfull lines
\setcounter{secnumdepth}{5}
% Make \paragraph and \subparagraph free-standing
\makeatletter
\ifx\paragraph\undefined\else
  \let\oldparagraph\paragraph
  \renewcommand{\paragraph}{
    \@ifstar
      \xxxParagraphStar
      \xxxParagraphNoStar
  }
  \newcommand{\xxxParagraphStar}[1]{\oldparagraph*{#1}\mbox{}}
  \newcommand{\xxxParagraphNoStar}[1]{\oldparagraph{#1}\mbox{}}
\fi
\ifx\subparagraph\undefined\else
  \let\oldsubparagraph\subparagraph
  \renewcommand{\subparagraph}{
    \@ifstar
      \xxxSubParagraphStar
      \xxxSubParagraphNoStar
  }
  \newcommand{\xxxSubParagraphStar}[1]{\oldsubparagraph*{#1}\mbox{}}
  \newcommand{\xxxSubParagraphNoStar}[1]{\oldsubparagraph{#1}\mbox{}}
\fi
\makeatother


\providecommand{\tightlist}{%
  \setlength{\itemsep}{0pt}\setlength{\parskip}{0pt}}\usepackage{longtable,booktabs,array}
\usepackage{calc} % for calculating minipage widths
% Correct order of tables after \paragraph or \subparagraph
\usepackage{etoolbox}
\makeatletter
\patchcmd\longtable{\par}{\if@noskipsec\mbox{}\fi\par}{}{}
\makeatother
% Allow footnotes in longtable head/foot
\IfFileExists{footnotehyper.sty}{\usepackage{footnotehyper}}{\usepackage{footnote}}
\makesavenoteenv{longtable}
\usepackage{graphicx}
\makeatletter
\def\maxwidth{\ifdim\Gin@nat@width>\linewidth\linewidth\else\Gin@nat@width\fi}
\def\maxheight{\ifdim\Gin@nat@height>\textheight\textheight\else\Gin@nat@height\fi}
\makeatother
% Scale images if necessary, so that they will not overflow the page
% margins by default, and it is still possible to overwrite the defaults
% using explicit options in \includegraphics[width, height, ...]{}
\setkeys{Gin}{width=\maxwidth,height=\maxheight,keepaspectratio}
% Set default figure placement to htbp
\makeatletter
\def\fps@figure{htbp}
\makeatother

\usepackage{indentfirst}
\usepackage{float}
\floatplacement{figure}{H}
\IfFileExists{plex-otf.sty}{
  %% Full TeXlive
  \usepackage[
    % math,
    RM={Scale=0.94},SS={Scale=0.94},SScon={Scale=0.94},TT={Scale=MatchLowercase,FakeStretch=0.9},
    DefaultFeatures={Ligatures=Common}
  ]{plex-otf}
  % \usepackage[Scale=MatchUppercase]{juliamono}
}{
  %% TinyTeX
  \usepackage{libertine}
}
\KOMAoption{captions}{tableheading}
\makeatletter
\@ifpackageloaded{caption}{}{\usepackage{caption}}
\AtBeginDocument{%
\ifdefined\contentsname
  \renewcommand*\contentsname{Содержание}
\else
  \newcommand\contentsname{Содержание}
\fi
\ifdefined\listfigurename
  \renewcommand*\listfigurename{Список иллюстраций}
\else
  \newcommand\listfigurename{Список иллюстраций}
\fi
\ifdefined\listtablename
  \renewcommand*\listtablename{Список таблиц}
\else
  \newcommand\listtablename{Список таблиц}
\fi
\ifdefined\figurename
  \renewcommand*\figurename{Рисунок}
\else
  \newcommand\figurename{Рисунок}
\fi
\ifdefined\tablename
  \renewcommand*\tablename{Таблица}
\else
  \newcommand\tablename{Таблица}
\fi
}
\@ifpackageloaded{float}{}{\usepackage{float}}
\floatstyle{ruled}
\@ifundefined{c@chapter}{\newfloat{codelisting}{h}{lop}}{\newfloat{codelisting}{h}{lop}[chapter]}
\floatname{codelisting}{Список}
\newcommand*\listoflistings{\listof{codelisting}{Листинги}}
\makeatother
\makeatletter
\makeatother
\makeatletter
\@ifpackageloaded{caption}{}{\usepackage{caption}}
\@ifpackageloaded{subcaption}{}{\usepackage{subcaption}}
\makeatother

\ifLuaTeX
\usepackage[bidi=basic]{babel}
\else
\usepackage[bidi=default]{babel}
\fi
\babelprovide[main,import]{russian}
\babelprovide[import]{english}
% get rid of language-specific shorthands (see #6817):
\let\LanguageShortHands\languageshorthands
\def\languageshorthands#1{}
\ifLuaTeX
  \usepackage{selnolig}  % disable illegal ligatures
\fi
\usepackage[style=gost-numeric,backend=biber,langhook=extras,autolang=other*]{biblatex}
\addbibresource{bib/cite.bib}
\usepackage{csquotes}
\usepackage{bookmark}

\IfFileExists{xurl.sty}{\usepackage{xurl}}{} % add URL line breaks if available
\urlstyle{same} % disable monospaced font for URLs
\hypersetup{
  pdftitle={Дискреционное разграничение прав в Linux},
  pdfauthor={Симонова Полина Игоревна (да, именно симонова)},
  pdflang={ru-RU},
  hidelinks,
  pdfcreator={LaTeX via pandoc}}


\title{Дискреционное разграничение прав в Linux}
\usepackage{etoolbox}
\makeatletter
\providecommand{\subtitle}[1]{% add subtitle to \maketitle
  \apptocmd{\@title}{\par {\large #1 \par}}{}{}
}
\makeatother
\subtitle{Основные атрибуты}
\author{Симонова Полина Игоревна (да, именно симонова)}
\date{}

\begin{document}
\maketitle

\renewcommand*\contentsname{Содержание}
{
\setcounter{tocdepth}{1}
\tableofcontents
}
\listoffigures
\listoftables

\setstretch{1.5}
\chapter{Цель
работы}\label{ux446ux435ux43bux44c-ux440ux430ux431ux43eux442ux44b}

Получение практических навыков работы в консоли с атрибутами файлов,
закрепление теоретических основ дискреционного разграничения доступа в
современных системах с открытым кодом на базе ОС Linux

\chapter{Задание}\label{ux437ux430ux434ux430ux43dux438ux435}

\begin{enumerate}
\def\labelenumi{\arabic{enumi}.}
\tightlist
\item
  Работа с атрибутами файлов
\item
  Заполнение таблицы \enquote{Установленные права и разрешённые
  действия} (см. табл. 2.1)
\item
  Заполнение таблицы \enquote{Минимальные права для совершения операций}
  (см. табл. 2.2)
\end{enumerate}

\chapter{Теоретическое
введение}\label{ux442ux435ux43eux440ux435ux442ux438ux447ux435ux441ux43aux43eux435-ux432ux432ux435ux434ux435ux43dux438ux435}

\textbf{Операционная система} --- это комплекс программ, предназначенных
для управления ресурсами компьютера и организации взаимодействия с
пользователем.

\textbf{Права доступа} определяют, какие действия конкретный
пользователь может или не может совершать с определенным файлами и
каталогами. С помощью разрешений можно создать надежную среду --- такую,
в которой никто не может поменять содержимое ваших документов или
повредить системные файлы.

\chapter{Выполнение лабораторной
работы}\label{ux432ux44bux43fux43eux43bux43dux435ux43dux438ux435-ux43bux430ux431ux43eux440ux430ux442ux43eux440ux43dux43eux439-ux440ux430ux431ux43eux442ux44b}

\section{Атрибуты
файлов}\label{ux430ux442ux440ux438ux431ux443ux442ux44b-ux444ux430ux439ux43bux43eux432}

\begin{enumerate}
\def\labelenumi{\arabic{enumi}.}
\tightlist
\item
  В операционной системе Rocky создаю нового пользователя guest через
  учетную запись администратора (рис. 1).
\end{enumerate}

\begin{figure}

{\centering \includegraphics[width=0.7\textwidth,height=\textheight]{image/1.PNG}

}

\caption{Добавление пользователя}

\end{figure}%

\begin{enumerate}
\def\labelenumi{\arabic{enumi}.}
\setcounter{enumi}{1}
\tightlist
\item
  Далее задаю пароль для созданной учетной записи (рис. 2).
\end{enumerate}

\begin{figure}

{\centering \includegraphics[width=0.7\textwidth,height=\textheight]{image/2.PNG}

}

\caption{Добавление пароля для пользователя}

\end{figure}%

\begin{enumerate}
\def\labelenumi{\arabic{enumi}.}
\setcounter{enumi}{2}
\tightlist
\item
  Сменяю пользователя в системе на только что созданного пользователя
  guest (рис. 3).
\end{enumerate}

\begin{figure}

{\centering \includegraphics[width=0.7\textwidth,height=\textheight]{image/3.PNG}

}

\caption{Вход через аккаунт пользователя}

\end{figure}%

\begin{enumerate}
\def\labelenumi{\arabic{enumi}.}
\setcounter{enumi}{3}
\tightlist
\item
  Определяю с помощью команды pwd, что я нахожусь в директории
  /home/guest/. Эта директория является домашней, ведь в приглашении
  командой строкой стоит значок \textasciitilde, указывающий, что я в
  домашней директории (рис. 4).
\end{enumerate}

\begin{figure}

{\centering \includegraphics[width=0.7\textwidth,height=\textheight]{image/4.PNG}

}

\caption{Текущая директория}

\end{figure}%

\begin{enumerate}
\def\labelenumi{\arabic{enumi}.}
\setcounter{enumi}{4}
\tightlist
\item
  Уточняю имя пользователя (рис. 5)
\end{enumerate}

\begin{figure}

{\centering \includegraphics[width=0.7\textwidth,height=\textheight]{image/5.PNG}

}

\caption{Информация об имени пользователе}

\end{figure}%

\begin{enumerate}
\def\labelenumi{\arabic{enumi}.}
\setcounter{enumi}{5}
\tightlist
\item
  В выводе команды groups информация только о названии группы, к которой
  относится пользователь. В выводе команды id можно найти больше
  информации: имя пользователя и имя группы, также коды имени
  пользователя и группы (рис. 6)
\end{enumerate}

\begin{figure}

{\centering \includegraphics[width=0.7\textwidth,height=\textheight]{image/6.PNG}

}

\caption{Информация о пользователе}

\end{figure}%

\begin{enumerate}
\def\labelenumi{\arabic{enumi}.}
\setcounter{enumi}{6}
\tightlist
\item
  Имя пользователя в приглашении командной строкой совпадает с именем
  пользователя, которое выводит команда whoami (рис. 7)
\end{enumerate}

\begin{figure}

{\centering \includegraphics[width=0.7\textwidth,height=\textheight]{image/7.PNG}

}

\caption{Сравнение информации об имени пользователя}

\end{figure}%

\begin{enumerate}
\def\labelenumi{\arabic{enumi}.}
\setcounter{enumi}{7}
\tightlist
\item
  Получаю информацию о пользователе с помощью команды
\end{enumerate}

\begin{verbatim}
cat /etc/passwd | grep guest
\end{verbatim}

В выводе получаю коды пользователя и группы, адрес домашней директории
(рис. 8).

\begin{figure}

{\centering \includegraphics[width=0.7\textwidth,height=\textheight]{image/8.PNG}

}

\caption{Просмотр файла passwd}

\end{figure}%

\begin{enumerate}
\def\labelenumi{\arabic{enumi}.}
\setcounter{enumi}{8}
\tightlist
\item
  Список поддиректорий директории home получилось получить с помощью
  команды ls -l, если мы добавим опцию -a, то сможем увидеть еще и
  директорию пользователя root. Права у директории:
\end{enumerate}

root: drwxr-xr-x,

pisimonova и guest: drwx------ (рис. 9).

\begin{figure}

{\centering \includegraphics[width=0.7\textwidth,height=\textheight]{image/9.PNG}

}

\caption{Просмотр содержимого директории}

\end{figure}%

\begin{enumerate}
\def\labelenumi{\arabic{enumi}.}
\setcounter{enumi}{9}
\tightlist
\item
  Пыталась проверить расширенные атрибуты директорий. Нет, их увидеть не
  удалось (рис. 10). Увидеть расширенные атрибуты других пользователей,
  тоже не удалось, для них даже вывода списка директорий не было.
\end{enumerate}

\begin{figure}

{\centering \includegraphics[width=0.7\textwidth,height=\textheight]{image/10.PNG}

}

\caption{Проверка расширенных атрибутов}

\end{figure}%

\begin{enumerate}
\def\labelenumi{\arabic{enumi}.}
\setcounter{enumi}{10}
\tightlist
\item
  Создаю поддиректорию dir1 для домашней директории. Расширенные
  атрибуты командой lsattr просмотреть у директории не удается, но
  атрибуты есть: drwxr-xr-x, их удалось просмотреть с помощью команды ls
  -l (рис. 11).
\end{enumerate}

\begin{figure}

{\centering \includegraphics[width=0.7\textwidth,height=\textheight]{image/11.PNG}

}

\caption{Создание поддиректории}

\end{figure}%

\begin{enumerate}
\def\labelenumi{\arabic{enumi}.}
\setcounter{enumi}{11}
\tightlist
\item
  Снимаю атрибуты командой chmod 000 dir1, при проверке с помощью
  команды ls -l видно, что теперь атрибуты действительно сняты (рис.
  12).
\end{enumerate}

\begin{figure}

{\centering \includegraphics[width=0.7\textwidth,height=\textheight]{image/12.PNG}

}

\caption{Снятие атрибутов с директории}

\end{figure}%

\begin{enumerate}
\def\labelenumi{\arabic{enumi}.}
\setcounter{enumi}{12}
\tightlist
\item
  Попытка создать файл в директории dir1. Выдает ошибку:
  \enquote{Отказано в доступе} (рис. 13).
\end{enumerate}

\begin{figure}

{\centering \includegraphics[width=0.7\textwidth,height=\textheight]{image/13.PNG}

}

\caption{Попытка создания файла}

\end{figure}%

Вернув права директории и использовав снова командy ls -l можно
убедиться, что файл не был создан (рис. 14).

\begin{figure}

{\centering \includegraphics[width=0.7\textwidth,height=\textheight]{image/14.PNG}

}

\caption{Проверка содержимого директории}

\end{figure}%

\section{Заполнение таблицы
2.1}\label{ux437ux430ux43fux43eux43bux43dux435ux43dux438ux435-ux442ux430ux431ux43bux438ux446ux44b-2.1}

\begin{longtable}[]{@{}
  >{\raggedright\arraybackslash}p{(\columnwidth - 18\tabcolsep) * \real{0.1000}}
  >{\raggedright\arraybackslash}p{(\columnwidth - 18\tabcolsep) * \real{0.1000}}
  >{\raggedright\arraybackslash}p{(\columnwidth - 18\tabcolsep) * \real{0.1000}}
  >{\raggedright\arraybackslash}p{(\columnwidth - 18\tabcolsep) * \real{0.1000}}
  >{\raggedright\arraybackslash}p{(\columnwidth - 18\tabcolsep) * \real{0.1000}}
  >{\raggedright\arraybackslash}p{(\columnwidth - 18\tabcolsep) * \real{0.1000}}
  >{\raggedright\arraybackslash}p{(\columnwidth - 18\tabcolsep) * \real{0.1000}}
  >{\raggedright\arraybackslash}p{(\columnwidth - 18\tabcolsep) * \real{0.1000}}
  >{\raggedright\arraybackslash}p{(\columnwidth - 18\tabcolsep) * \real{0.1000}}
  >{\raggedright\arraybackslash}p{(\columnwidth - 18\tabcolsep) * \real{0.1000}}@{}}
\toprule\noalign{}
\endhead
\bottomrule\noalign{}
\endlastfoot
Права директории & Права файла & Создание файла & Удаление файла &
Запись в файл & Чтение файла & Смена директории & Просмотр файлов в
директории & Переимено- вание файла & Смена атрибутов файла \\
d(000) & (000) & - & - & - & - & - & - & - & - \\
d(000) & (100) & - & - & - & - & - & - & - & - \\
d(000) & (200) & - & - & - & - & - & - & - & - \\
d(000) & (300) & - & - & - & - & - & - & - & - \\
d(000) & (400) & - & - & - & - & - & - & - & - \\
d(000) & (500) & - & - & - & - & - & - & - & - \\
d(000) & (600) & - & - & - & - & - & - & - & - \\
d(000) & (700) & - & - & - & - & - & - & - & - \\
d(100) & (000) & - & - & - & - & + & - & - & + \\
d(100) & (100) & - & - & - & - & + & - & - & + \\
d(100) & (200) & - & - & + & - & + & - & - & + \\
d(100) & (300) & - & - & + & - & + & - & - & + \\
d(100) & (400) & - & - & - & + & + & - & - & + \\
d(100) & (500) & - & - & - & + & + & - & - & + \\
d(100) & (600) & - & - & + & + & + & - & - & + \\
d(100) & (700) & - & - & + & + & + & - & - & + \\
d(200) & (000) & - & - & - & - & - & - & - & - \\
d(200) & (100) & - & - & - & - & - & - & - & - \\
d(200) & (200) & - & - & - & - & - & - & - & - \\
d(200) & (300) & - & - & - & - & - & - & - & - \\
d(200) & (400) & - & - & - & - & - & - & - & - \\
d(200) & (500) & - & - & - & - & - & - & - & - \\
d(200) & (600) & - & - & - & - & - & - & - & - \\
d(200) & (700) & - & - & - & - & - & - & - & - \\
d(300) & (000) & + & + & - & - & + & - & + & + \\
d(300) & (100) & + & + & - & - & + & - & + & + \\
d(300) & (200) & + & + & + & - & + & - & + & + \\
d(300) & (300) & + & + & + & - & + & - & + & + \\
d(300) & (400) & + & + & - & + & + & - & + & + \\
d(300) & (500) & + & + & - & + & + & - & + & + \\
d(300) & (600) & + & + & + & + & + & - & + & + \\
d(300) & (700) & + & + & + & + & + & - & + & + \\
d(400) & (000) & - & - & - & - & - & + & - & - \\
d(400) & (100) & - & - & - & - & - & + & - & - \\
d(400) & (200) & - & - & - & - & - & + & - & - \\
d(400) & (300) & - & - & - & - & - & + & - & - \\
d(400) & (400) & - & - & - & - & - & + & - & - \\
d(400) & (500) & - & - & - & - & - & + & - & - \\
d(400) & (600) & - & - & - & - & - & + & - & - \\
d(400) & (700) & - & - & - & - & - & + & - & - \\
d(500) & (000) & - & - & - & - & + & + & - & + \\
d(500) & (100) & - & - & - & - & + & + & - & + \\
d(500) & (200) & - & - & + & - & + & + & - & + \\
d(500) & (300) & - & - & + & - & + & + & - & + \\
d(500) & (400) & - & - & - & + & + & + & - & + \\
d(500) & (500) & - & - & - & + & + & + & - & + \\
d(500) & (600) & - & - & + & + & + & + & - & + \\
d(500) & (700) & - & - & + & + & + & + & - & + \\
d(600) & (000) & - & - & - & - & - & + & - & - \\
d(600) & (100) & - & - & - & - & - & + & - & - \\
d(600) & (200) & - & - & - & - & - & + & - & - \\
d(600) & (300) & - & - & - & - & - & + & - & - \\
d(600) & (400) & - & - & - & - & - & + & - & - \\
d(600) & (500) & - & - & - & - & - & + & - & - \\
d(600) & (600) & - & - & - & - & - & + & - & - \\
d(600) & (700) & - & - & - & - & - & + & - & - \\
d(700) & (000) & + & + & - & - & + & + & + & + \\
d(700) & (100) & + & + & - & - & + & + & + & + \\
d(700) & (200) & + & + & + & - & + & + & + & + \\
d(700) & (300) & + & + & + & - & + & + & + & + \\
d(700) & (400) & + & + & - & + & + & + & + & + \\
d(700) & (500) & + & + & - & + & + & + & + & + \\
d(700) & (600) & + & + & + & + & + & + & + & + \\
d(700) & (700) & + & + & + & + & + & + & + & + \\
\end{longtable}

Таблица 2.1 «Установленные права и разрешённые действия»

Пример заполнения таблицы 2.1 (рис. 15).

\begin{figure}

{\centering \includegraphics[width=0.7\textwidth,height=\textheight]{image/15.PNG}

}

\caption{Изменение прав директории и файла}

\end{figure}%

\section{Заполнение таблицы
2.2}\label{ux437ux430ux43fux43eux43bux43dux435ux43dux438ux435-ux442ux430ux431ux43bux438ux446ux44b-2.2}

\begin{longtable}[]{@{}
  >{\raggedright\arraybackslash}p{(\columnwidth - 8\tabcolsep) * \real{0.2000}}
  >{\raggedright\arraybackslash}p{(\columnwidth - 8\tabcolsep) * \real{0.2000}}
  >{\raggedright\arraybackslash}p{(\columnwidth - 8\tabcolsep) * \real{0.2000}}
  >{\raggedright\arraybackslash}p{(\columnwidth - 8\tabcolsep) * \real{0.2000}}
  >{\raggedright\arraybackslash}p{(\columnwidth - 8\tabcolsep) * \real{0.2000}}@{}}
\toprule\noalign{}
\endhead
\bottomrule\noalign{}
\endlastfoot
Операция & & Минимальные права на директорию & & Минимальные права на
файл \\
Создание файла & & d(300) & & - \\
Удаление файла & & d(300) & & - \\
Чтение файла & & d(100) & & (400) \\
Запись в файл & & d(100) & & (200) \\
Переименование файла & & d(300) & & (000) \\
Создание поддиректории & & d(300) & & - \\
Удаление поддиректории & & d(300) & & - \\
\end{longtable}

Таблица 2.2 \enquote{Минимальные права для совершения операций}

\chapter{Выводы}\label{ux432ux44bux432ux43eux434ux44b}

Были получены практические навыки работы в консоли с атрибутами файлов,
закреплены теоретические основы дискреционного разграничения доступа в
современных системах с открытым кодом на базе ОС Linux.

\chapter{Список
литературы}\label{ux441ux43fux438ux441ux43eux43a-ux43bux438ux442ux435ux440ux430ux442ux443ux440ux44b}

{[}1{]} Операционные системы:
https://blog.skillfactory.ru/glossary/operaczionnaya-sistema/


\printbibliography



\end{document}
